\section{Actual nonflat finite-map principal degrees on the original curve charts}

\begin{theorem}[Whole original affine finite-map principal degree]
\label{thm:fc-scheme-finite-principal-generic-rank}
\leanok
\lean{EnvelopeProperty.SchemeFinitePrincipalDegrees.nonempty_preimage, EnvelopeProperty.SchemeFinitePrincipalDegrees.sections_krullDimLE, EnvelopeProperty.SchemeFinitePrincipalDegrees.affinePrincipal_length_eq_genericRank}
\uses{thm:fc-finite-algebra-principal-degree}
For any original finite dominant map of integral schemes to an
original locally Noetherian scheme of dimension at most one, and
any original nonempty affine open, the actual whole inverse-chart
section algebra quotient by any nonzero original chart section has
original base-ring length equal to its true generic rank times the
native order of that SAME original section. No flatness or
normality is assumed.
\end{theorem}
\begin{proof}
Actual dominance sends original generic points to one another,
so the whole inverse chart is genuinely nonempty. Original
integrality makes both actual section rings domains. Native finite
morphism geometry gives the original whole section map finite,
and actual dominance gives injectivity. Original affine scheme
geometry derives Noetherianity and the dimension bound of the
SAME original chart ring. The genuine finite-algebra principal
length theorem therefore applies with no algebraic comparison
inputs.
\end{proof}

\begin{theorem}[The multiplier is the same original function-field degree]
\label{thm:fc-scheme-finite-principal-functionfield-degree}
\leanok
\lean{EnvelopeProperty.SchemeFinitePrincipalDegrees.affinePrincipal_length_eq_functionFieldDegree}
\uses{thm:fc-scheme-finite-principal-generic-rank}
In the preceding true principal degree formula, the multiplier is
the degree of the SAME original native function-field map of the
actual original dominant morphism. All original fraction-field
structures and their coefficient squares are internally derived.
\end{theorem}
\begin{proof}
The actual original nonempty affine and inverse-affine charts have
their actual original whole scheme function fields as fraction
fields of the SAME original section rings. The native original
dominant function-field map retains every original generic germ
under the actual whole section map. This internally proves both
literal coefficient towers; the original fraction-field rank
formula then identifies the multiplier with the degree of this
SAME original function-field extension. No chosen field embedding,
rank, length or matching certificate is supplied.
\end{proof}
